% =============================================================================
% Paper 2 R1 SEVERE 2 — Linearity Robustness of VIXTWN-to-VIX Ratio
% Suggested placement: paper/taiwan-vt/body.tex
%   (a) New \subsection at end of §3 (data section), OR end of §4
%       VT strategies before §5 high-frequency section.
%   (b) Plus a one-sentence footnote attached to line 120 disclosing
%       regime dependence.
%
% Word count: 196 (target ~150-200 EN words; PBFJ-tier concise).
% Data source: experiments/paper2_R1_linear_scaling_fix/results.json
% =============================================================================

\subsection{Linearity Robustness of the VIXTWN-to-VIX Ratio}\label{sec:vixtwn_ratio_linearity}

The calibration $K = 12/1.39 = 8.63$ used in the VIX-proxy VT strategy
assumes the VIXTWN-to-VIX ratio is approximately stable across volatility
regimes. A referee may legitimately ask whether this linearity assumption
breaks in tail-volatility episodes, in which case $K$ would systematically
misallocate during the precise periods VT strategies aim to manage.
Table~\ref{tab:ratio_linearity} reports the ratio conditional on
expanding-window VIX quantile buckets (no full-sample lookahead) and on a
tail bucket defined as $|\Delta \log \text{VIX}_t| > 2\,\sigma_{t-1}$, using
the TAIFEX-official VIXTWN series (K1098 parse, $n=3{,}075$ post-warmup,
2008--2021).

The ratio declines monotonically from $1.018$ at low VIX (Q1) to $0.824$
at high VIX (Q4), a $19\%$ relative drop; the tail bucket averages $0.877$.
Each bucket's $95\%$ stationary-block bootstrap confidence interval
($B=500$, mean block $=21$ days, seed $=42$) excludes $1.39$. A
Kolmogorov-Smirnov two-sample test rejects equality of the tail and
non-tail ratio distributions ($D=0.274$, $p < 0.001$). The static-ratio
assumption therefore does not hold strictly, but two mitigating points
preserve the substantive conclusions of Section~\ref{sec:vt}. First, the
Sharpe ratio of the $8.63/\text{VIX}$ strategy is invariant to $K$
(Section~\ref{sec:vt}, line $\sim$313); $K$ affects only the risk level
and hence maximum drawdown. Second, the recent (Dec 2025--Apr 2026,
$n=76$) ratio of $1.393$ used in the calibration reflects a structurally
higher-volatility regime in Taiwan; our recommendation is to interpret
$K=8.63$ as a recent-window practical calibration and to disclose the
regime dependence rather than to revise the calibration.

\begin{table}[!htbp]
\centering
\caption{VIXTWN-to-VIX Ratio Conditional on VIX Regime, 2008--2021}
\label{tab:ratio_linearity}
\small
\begin{tabular}{lrrrrcc}
\toprule
Bucket & $n$ & Mean & Median & SD & 95\% CI & Contains 1.39? \\
\midrule
Q1 (VIX $\le$ Q25) & 1{,}533 & 1.018 & 1.018 & 0.141 & $[0.987, 1.046]$ & No \\
Q2 (Q25--Q50)      &   634 & 1.007 & 0.990 & 0.180 & $[0.955, 1.054]$ & No \\
Q3 (Q50--Q75)      &   574 & 0.947 & 0.919 & 0.176 & $[0.897, 0.995]$ & No \\
Q4 (VIX $>$ Q75)   &   334 & 0.824 & 0.824 & 0.133 & $[0.785, 0.862]$ & No \\
Tail ($2\sigma$)   &   183 & 0.877 & 0.877 & 0.202 & $[0.832, 0.919]$ & No \\
\midrule
Overall (long sample) & 3{,}075 & 0.981 & 0.985 & 0.180 & $[0.953, 1.009]$ & No \\
Recent (Dec 2025--Apr 2026) &  76 & 1.393 &     -- &     -- & --- & --- \\
\bottomrule
\end{tabular}
% source: experiments/paper2_R1_linear_scaling_fix/results.json
%   amplification_per_quantile.{Q1,Q2,Q3,Q4,Tail}.{mean_ratio,median_ratio,std_ratio,n}
%   bootstrap_ci.per_bucket.{Q1,Q2,Q3,Q4,Tail}.{ci_low,ci_high,contains_1_39}
%   bootstrap_ci.overall.{ci_low,ci_high,contains_1_39}
%   Recent-window row sourced from experiments/k1181/k1181_results.json
%     .vixtwn_ratio.recent_Dec2025_Apr2026.{mean,n}
\smallskip
\\
\small \textit{Notes:} Quantile thresholds are computed in an
expanding-window manner using observations strictly prior to day $t$;
the first 252 days are dropped as warmup. The tail bucket uses
$|\Delta \log \text{VIX}_t| > 2\,\sigma_{t-1}$ where $\sigma_{t-1}$ is the
expanding standard deviation of log-VIX changes through $t-1$. Confidence
intervals are computed via stationary block bootstrap
\citep{politis1994stationary} with $B=500$ replications and mean block
length $21$ trading days. The long-history sample uses the TAIFEX-official
VIXTWN series; the recent window uses K1181's daily snapshot. The decline
in the ratio at high VIX reflects under-reaction of VIXTWN to U.S.\
vol shocks rather than over-reaction in calm periods; mechanistically
this is consistent with structural-break evidence around the 2020 launch
of an active VIXTWN futures market.
\end{table}

% =============================================================================
% Optional footnote variant (if main-thread chooses footnote-only treatment
% instead of new subsection). Attach to body.tex line 120 after
% "...higher baseline volatility of the Taiwanese market relative to the
% S&P 500.\footnote{...}"
% =============================================================================
%
%   \footnote{The single figure of $1.393$ reflects the recent
%   Dec 2025--Apr 2026 window ($n=76$, K1181). Over the longer 2008--2021
%   TAIFEX-official VIXTWN sample ($n=3{,}075$), the ratio averages
%   $0.981$ and varies systematically with VIX regime, ranging from
%   $1.018$ at low VIX to $0.824$ at high VIX (each 95\% bootstrap CI
%   excludes 1.39; see
%   Section~\ref{sec:vixtwn_ratio_linearity}/Table~\ref{tab:ratio_linearity}).
%   Because $\text{Sharpe}(K \cdot r) = \text{Sharpe}(r)$, this regime
%   dependence affects the calibration of $K$ and hence the strategy's
%   target risk level, but not its Sharpe ranking.}
